% SPDX-License-Identifier: Apache-2.0
\section{Two Clocks, the Slow One First}
\label{sec:x-slowfirst}

The ports file lists a netlist's clocks in the order the lowering meets
them. \code{SlowFirst} holds the two counters of
Section~\ref{sec:x-twoclock} with the slow one first, so its ports file
names \code{slow} before \code{clk}. A port associated with another
clock domain includes an explicit annotation, \code{@slow}, while
ports in the default clock domain remain unannotated.

The testbench generator read an unmarked port as on the first clock the
file lists. Here that put \code{go} and \code{a}, which are on
\code{clk}, on \code{slow}'s edges: \code{go} was taken from the trace
at the next slow edge, so the netlist was told it for three cycles
where the run had it for one, and \code{a} was checked a slow period
later. Both simulators then reported \code{a} wrong from tick 17 on,
and the netlist was not at fault. An unmarked port is on \code{clk}
now whenever the netlist names \code{clk}, and on the first clock only
when it does not (issue 888).

\code{go} is high for one cycle, at tick 18, between two edges of
\code{slow}, which is where reading it at the wrong clock shows. With
the fix both simulators agree with the run.

\begin{figure}[htbp]
\centering
\input{slowfirst_timing.tex}
\caption{\code{go} high for one cycle of \code{clk} at tick 18, and
\code{a} counting one; \code{b} counts on the slow clock.}
\label{fig:x-slowfirst}
\end{figure}

\lstinputlisting[rangebeginprefix=//\ begin\{,rangebeginsuffix=\},%
rangeendprefix=//\ end\{,rangeendsuffix=\},includerangemarker=false,%
linerange=units-units,caption={\code{lib/examples/ex\_slowfirst.rs}:
the two counters and the parent that holds the slow one
first.},label={lst:x-slowfirst}]{lib/examples/ex_slowfirst.rs}

\lstinputlisting[style=ascii,caption={What \code{ex\_slowfirst} printed
when the build ran it: the two counts, then the ports file, with
\code{slow} first and no mark on \code{go} or \code{a}.},label={lst:x-slowfirst-out}]{out_slowfirst.txt}
